% SGA 3, 第 VI_B 篇, 第 1 轮英文重构底本的中译.
% 范围: 注 7.6.0 与命题 7.6 全部, 包括特例段落、完整证明及原编者注 76--77.
% 权威底本: Polo--Gille Exp6B-13oct24.pdf, 组件第 56--57 页
% (合卷第 397--398 页; 重版页眉第 386--387 页; 原版边栏页码第 370 页).
% 不含组件第 56 页的前一引理 7.5 及第 57 页的后一注 7.6.1.
% OCR 与 Reinhold 英文 Markdown 仅为对照材料. 本组件无图.

\medskip
\phantomsection\label{rem:sga3-VIB-7.6.0}%
\noindent\textbf{注 7.6.0.---}%
\footnote{\label{note:sga3-VIB-76}编者注: 此定义为新增内容;
在原文中, 它包含于\hyperref[prop:sga3-VIB-7.6]{命题 7.6} 的陈述中.}
\emph{设 \(S\) 是一个概形, \(H\) 是一个 \(S\)-群,
\(f:X\to H\) 是一个 \(S\)-概形态射.
由 \(f\) 的像生成的 \(H\) 的群子预层, 记作}
\[
  \bigl\langle\operatorname{Im}f\bigr\rangle,
\]
\emph{是 \(H\) 的如下群子函子: 它将每个 \(S\)-概形 \(S'\)
映为由 \(f(X(S'))\) 生成的 \(H(S')\) 的子群.
由于此子群的每个元素都可写作有限乘积}
\phantomsection\label{form:sga3-VIB-7.6.0-finite-product}%
\[
  f(x_1)^{\varepsilon_1}\cdots f(x_n)^{\varepsilon_n},
  \qquad
  n\in\mathbb{N}^{*},\quad
  x_i\in X(S'),\quad
  \varepsilon_i=\pm1,
\]
\emph{所以, 若置}
\[
  X^1=X\amalg X
\]
\emph{并如 \hyperref[concl:sga3-VIB-7.1-ii]{7.1} 那样定义态射
\(f^n:X^n\to H\), 则 \(\langle\operatorname{Im}f\rangle\)
恰好是态射}
\phantomsection\label{form:sga3-VIB-7.6.0-finfty}%
\[
  f^\infty:X^\infty\longrightarrow H,
\]
\emph{的像预层, 其中}
\[
  X^\infty=\coprod_{n\geq1}X^n
\]
\emph{而 \(f^\infty:X^\infty\to H\) 是在每个 \(X^n\)
上的限制为 \(f^n\) 的 \(S\)-态射.}

\medskip
\phantomsection\label{prop:sga3-VIB-7.6}%
\noindent\textbf{命题 7.6.---}%
\footnote{\label{note:sga3-VIB-77}编者注: 原文在特例的假设下陈述此结果,
但在 \SGARefLink{prop:sga3-VIB-10.12}{10.12} 的证明中隐含地使用了更一般的形式;
因此相应地改写了陈述.}
\emph{设 \(S\) 是一个概形, \(X\) 是一个在 \(S\) 上平坦且局部有限呈示的
\(S\)-概形, \(H\) 是一个在 \(S\) 上局部有限呈示且纤维约化的
\(S\)-群, \(f:X\to H\) 是一个 \(S\)-概形态射,
\(f^\infty:X^\infty\to H\) 是以上引入的 \(S\)-态射.
以下条件等价:}
\begin{enumerate}
\item[\textup{\SGAInlineRefTarget{prop:sga3-VIB-7.6:item:i}(i)}]\phantomsection\label{cond:sga3-VIB-7.6-i}
  \emph{\(H\) 表出与预层 \(\langle\operatorname{Im}f\rangle\)
  相伴的关于 \((\mathrm{fppf})\) 拓扑的 \(S\)-层.}

\item[\textup{\SGAInlineRefTarget{prop:sga3-VIB-7.6:item:ii}(ii)}]\phantomsection\label{cond:sga3-VIB-7.6-ii}
  \emph{\(f^\infty\) 是 \((\mathrm{fppf})\) 拓扑的覆盖.}

\item[\textup{\SGAInlineRefTarget{prop:sga3-VIB-7.6:item:iii}(iii)}]\phantomsection\label{cond:sga3-VIB-7.6-iii}
  \emph{\(f^\infty\) 是满射, 即}
  \phantomsection\label{form:sga3-VIB-7.6-union}%
  \[
    H=\bigcup_{n\geq1}f^n(X^n).
  \]
\end{enumerate}
\emph{此外, 若 \(H\) 拟紧, 则这些条件还等价于以下条件:}
\begin{enumerate}
\item[\textup{(iv)}]\phantomsection\label{cond:sga3-VIB-7.6-iv}
  \emph{存在一个整数 \(n\), 使得 \(f^n:X^n\to H\)
  是 \((\mathrm{fppf})\) 拓扑的覆盖, 因而为满射.}
\end{enumerate}

\emph{这特别适用于以下情形:
\(X\) 是局部有限型的几何约化 \(k\)-概形,
\(G\) 是局部有限型的 \(k\)-群,
\(\varphi:X\to G\) 是 \(k\)-概形态射,
\(H=\Gamma_G(\varphi)\), 而 \(f:X\to H\)
是 \(\varphi\) 诱导的态射.}

\emph{证明.}
\textup{\SGARefLink{prop:sga3-VIB-7.6:item:i}{(i)}} 中考虑的层,
是 \(X^\infty\) 在 \(f^\infty\) 下的像层.
所以, 由 \hyperref[sga3:iv:proposition:4:4:3]{IV, 4.4.3},
\(H\) 表出此层等价于 \(f^\infty\) 是覆盖;
这又蕴含 \(f^\infty\) 是满射.
反之, 假设 \(f^\infty\) 为满射.
我们要证明它此时是 \((\mathrm{fppf})\) 拓扑的覆盖.

设 \(s\in S\), \(\eta\) 是纤维 \(H_s\) 的一个极大点,
\(x\in X^\infty_s\) 满足
\[
  f^\infty(x)=\eta;
\]
这样的 \(x\) 存在, 因为 \(f^\infty\) 是满射.
由于纤维 \(H_s\) 约化, \(\mathcal{O}_{H_s,\eta}\) 是域,
所以 \(f^\infty_s\) 在 \(x\) 处平坦.
由于 \(X^\infty\) 与 \(H\) 在 \(S\) 上局部有限呈示,
且 \(X^\infty\) 在 \(S\) 上平坦,
平坦性的逐纤维判别法 (EGA~IV\(_3\), 11.3.10)
表明 \(f^\infty\) 在 \(x\) 处平坦.
所以, 由\hyperref[lem:sga3-VIB-7.5]{引理 7.5}, 态射
\phantomsection\label{form:sga3-VIB-7.6-finfty-product}%
\[
  \mu\circ
  \bigl(f^\infty\times_S f^\infty\bigr):
  X^\infty\times_S X^\infty\longrightarrow H
\]
是 \((\mathrm{fppf})\) 拓扑的覆盖.

现在, \(X^n\times_S X^m\) 典范同构于 \(X^{n+m}\),
且在此同构下,
\[
  \mu\circ\bigl(f^n\times_S f^m\bigr)
\]
对应于 \(f^{n+m}\). 所以显然
\[
  \mu\circ
  \bigl(f^\infty\times_S f^\infty\bigr)
\]
通过 \(f^\infty\) 分解, 因而 \(f^\infty\)
是 \((\mathrm{fppf})\) 拓扑的覆盖.
这就证明 \(\textup{(iii)}\Rightarrow\textup{(ii)}\),
从而得到条件 \textup{\SGARefLink{prop:sga3-VIB-7.6:item:i}{(i)}},
\textup{\SGARefLink{prop:sga3-VIB-7.6:item:ii}{(ii)}} 与
\textup{\SGARefLink{prop:sga3-VIB-7.6:item:iii}{(iii)}} 的等价性.

此外, 由于态射 \(X\to S\) 与 \(H\to S\) 局部有限呈示,
\(f:X\to H\) 也局部有限呈示 (EGA~IV\(_1\), 1.4.3\textup{(v)}).
乘法
\[
  \mu:H\times_S H\longrightarrow H
\]
也局部有限呈示 (\hyperref[sga3:VIA:0.1]{VI\(_A\), 0.1}),
从而每个 \(f^n:X^n\to H\) 都局部有限呈示.
所以, 由 EGA~IV\(_1\), 1.9.5\textup{(viii)}, 集合
\[
  f^n(X^n)
\]
是 \(H\) 的归纳可构造子集.

再假设 \(H\) 拟紧; 则由于 \(H\to S\) 为满射, \(S\) 也拟紧.
由 EGA~IV\(_1\), 1.9.9, 存在一个整数 \(p\), 使得
\[
  f^p(X^p)=H.
\]
如前, 利用 \(H\) 的纤维约化这一事实,
结合\hyperref[lem:sga3-VIB-7.5]{引理 7.5}, 可知态射
\phantomsection\label{form:sga3-VIB-7.6-f2p-covering}%
\[
  \mu\circ\bigl(f^p\times_S f^p\bigr):
  X^p\times_S X^p\longrightarrow H
\]
是 \((\mathrm{fppf})\) 拓扑的覆盖.
由于此态射是
\[
  f^{2p}:X^{2p}\longrightarrow H,
\]
\hyperref[prop:sga3-VIB-7.6]{命题 7.6} 的证明完成.

% 中文翻译检查点：本组件全部正文、证明、脚注及源注已译；源183行。
% 机械核对：行内数学97处、陈列数学15处及全部引用键守恒（数学中的自然语言已译）。
% 状态：完成源文覆盖和逐段自校；尚未经独立合卷编译与页面审查。
